\documentclass[10pt,twocolumn]{article}
\usepackage[margin=0.58in]{geometry}
\usepackage{amsmath,amssymb,mathtools,bm}
\usepackage{microtype}
\usepackage{enumitem}
\usepackage{xcolor}
\usepackage{hyperref}
\usepackage{listings}
\hypersetup{colorlinks=true,linkcolor=blue!50!black,urlcolor=blue!50!black}
\setlength{\parindent}{0pt}
\setlength{\parskip}{2.2pt}
\setlist{nosep,leftmargin=*}
\lstset{
  language=Python,
  basicstyle=\ttfamily\scriptsize,
  columns=fullflexible,
  keepspaces=true,
  breaklines=true,
  breakatwhitespace=false,
  frame=single,
  framerule=0.25pt,
  xleftmargin=0pt,
  xrightmargin=0pt,
  aboveskip=3pt,
  belowskip=3pt,
  showstringspaces=false
}

\newcommand{\ket}[1]{\left|#1\right\rangle}
\newcommand{\bra}[1]{\left\langle#1\right|}
\newcommand{\braket}[2]{\left\langle#1\middle|#2\right\rangle}
\newcommand{\Acal}{\mathcal A}
\newcommand{\Gcal}{\mathcal G}
\newcommand{\Pcal}{\mathcal P}
\newcommand{\Qcal}{\mathcal Q}
\newcommand{\Rcal}{\mathcal R}
\newcommand{\Wcal}{\mathcal W}
\newcommand{\Sadapt}{S_{\rm ADAPT}}
\newcommand{\Ssnake}{S_{\rm SNAKE}}
\newcommand{\Salg}{S_{\rm alg}}
\newcommand{\Ngrad}{N_{\rm grad}}
\newcommand{\Nmetric}{N_{\rm metric}}
\newcommand{\NH}{N_H}
\newcommand{\NHouter}{N_{H,\rm outer}}
\newcommand{\NHrefit}{N_{H,\rm refit}}
\newcommand{\oneGeo}{\mathbf 1_{\rm Geo}}
\newcommand{\boxline}{\vspace{2pt}\hrule\vspace{3pt}}
\newenvironment{primerbox}[1]{%
  \boxline\begingroup\raggedright\textbf{#1.}\ }{\par\endgroup\boxline}

\begin{document}
\title{Paper I Shot Proxy: Reporting Primer}
\author{Paper-I support note}
\date{2026-06-11}
\maketitle
\vspace{-1.1em}

\begin{primerbox}{Scope}
This note describes the \(S\) value reported in the Paper I tables.  It does
not use the Math.pdf candidate-local shot-cost expression.  The reported \(S\)
is a normalized estimator/probe-work ledger through the displayed accepted
prefix.  Its units are counted controller probe records and Hamiltonian
objective evaluations, not calibrated physical circuit shots.
\end{primerbox}

\section*{1. Displayed Prefix}

Let \(k=0,1,\ldots\) index accepted adaptive iterations.  The displayed table
row uses an accepted prefix \(k_\star\): the first prefix whose energy supports
the displayed accuracy resolution.  Every reported sum below stops at
\(k_\star\).

For a non-SNAKE adaptive comparator, let \(\Acal_k\) be the live candidate set
scored before admission \(k\), and let \(n_k=|\Acal_k|\).  If the comparator is
geometric, the centered candidate tangents
\[
  \tau_r^{(k)}
  =
  (I-\ket{\psi_k}\bra{\psi_k})(-iA_r\ket{\psi_k})
\]
define
\[
  G_{rs}^{(k)}
  =
  \operatorname{Re}\braket{\tau_r^{(k)}}{\tau_s^{(k)}},
  \qquad r,s\in\Acal_k.
\]
The symmetric metric triangle contributes \(n_k(n_k+1)/2\).  If \(f_k\)
denotes the Hamiltonian objective evaluations charged after admission \(k\),
then
\[
\boxed{
\Sadapt
=
\sum_{k=0}^{k_\star}
\left(
n_k+\oneGeo\frac{n_k(n_k+1)}{2}+f_k
\right).
}
\]

\section*{2. From A Record To A Countable Probe}

SNAKE does not report a separate shot proxy for every low-level Pauli term.
The reporting path starts from candidate-position records.  A record is
\[
  r=(a,p),
\]
where \(a\) is a candidate generator and \(p\) is the proposed insertion
position in the current ansatz prefix.  The record carries an ansatz generator
\[
  A_r=\sum_{\nu} c_{r\nu}P_{r\nu},
  \qquad
  P_{r\nu}\in\{e,x,y,z\}^{\otimes n}.
\]
For a Pauli key \(u\in\{e,x,y,z\}^{\otimes n}\), write \(w(u)\) for its
non-identity weight.  A basis key \(v\) covers \(u\), written \(u\preceq v\),
when
\[
  \forall q,\qquad u_q=e\quad\hbox{or}\quad u_q=v_q .
\]
Two keys merge when every qubit is compatible:
\[
  (u\vee v)_q=
  \begin{cases}
  v_q, & u_q=e,\\
  u_q, & v_q=e,\\
  u_q, & u_q=v_q,
  \end{cases}
\]
and the merge is undefined when some non-identity axes disagree.

The reporting group-key map is the same QWC compression used by
\texttt{measurement\_group\_keys\_for\_term}:
\[
  \Gamma(r)
  =
  {\rm Compress}_{\rm QWC}
  \{\,P_{r\nu}: |c_{r\nu}|>0,\; w(P_{r\nu})>0\,\}.
\]
A selector record is measurement-bearing exactly when
\[
  \Gamma(r)\ne\varnothing .
\]
The strict Paper I SNAKE reconstruction reads
\texttt{records\_with\_group\_keys}.  Thus the count for a phase is the number
of measurement-bearing records, not \(\sum_r|\Gamma(r)|\).

\[
\boxed{
  N_{\rm phase}(k)
  =
  |\{r\in\Rcal_{\rm phase}(k):\Gamma(r)\ne\varnothing\}|.
}
\]

\section*{3. How \(P_1,P_2,P_3\) Are Found}

At prefix \(k\), let \(\Omega_k\) be the available candidate-position surface.
The controller scores records on this surface and moves only selected records
to later phases.  Write
\[
  {\rm Short}_{j,k}(\Rcal;s,\tau,B)
\]
for the deterministic controller shortlist: sort records by score \(s\), keep
records above the phase threshold/frontier rule, and cap by the phase budget
\(B\).  This symbol names the table-support selection operation; the \(S\)
column counts the records it causes to be probed.

The Phase-I record surface is
\[
  \begin{aligned}
  \Rcal_1(k)
  =
  \{(a,p)\in\Omega_k:\;&
  (a,p)\hbox{ survives}\\
  &\hbox{the Phase-0 pilot screen}\}.
  \end{aligned}
\]
The Paper I Phase-I probe set is the measurement-bearing part of that surface:
\[
  P_1(k)=\{r\in\Rcal_1(k):\Gamma(r)\ne\varnothing\}.
\]
The Phase-I shortlist is
\[
  \Qcal_1(k)
  =
  {\rm Short}_{1,k}(\Rcal_1(k);s_1,\tau_1,B_1).
\]

Phase II reranks the Phase-I shortlist with the full geometric score.  Its
reported probe surface is
\[
  \Rcal_2(k)
  =
  {\rm FullGeom}_{2,k}\bigl(\Qcal_1(k)\bigr),
\]
where each retained record is rebuilt with the current gradient, tangent norm,
hardware cost, measurement-stat payload, and Phase-II score.  The reported
Phase-II set is
\[
  P_2(k)=\{r\in\Rcal_2(k):\Gamma(r)\ne\varnothing\}.
\]
The Phase-II shortlist is
\[
  \Qcal_2(k)
  =
  {\rm Short}_{2,k}(\Rcal_2(k);s_2,\tau_2,B_2).
\]

Phase III reranks the Phase-II shortlist with the reduced Schur-window
geometry.  The reported Phase-III probe surface is
\[
  \Rcal_3(k)=\Qcal_2(k),
  \qquad
  P_3(k)=\{r\in\Rcal_3(k):\Gamma(r)\ne\varnothing\}.
\]
The final admission set is selected after this rerank:
\[
  \Qcal_3(k)
  =
  {\rm Short}_{3,k}(\Rcal_3(k);s_3,\tau_3,B_3).
\]
The table \(S\) charges the probes used to form \(\Rcal_1,\Rcal_2,\Rcal_3\).
It does not replace those counts by the size of the final admitted batch.
Since Phase II is built from the Phase-I shortlist and Phase III is built from
the Phase-II shortlist, the same-prefix record counts obey
\[
  |P_3(k)|\le |P_2(k)|\le |P_1(k)|.
\]

\begin{primerbox}{Phase-0 gradient screen}
The implementation also has a Phase-0 pilot/selected-logical gradient screen.
The table reconstruction puts it in the gradient bin with Phase I:
\[
  \Ngrad
  =
  \sum_{k=0}^{k_\star}
  \left(N_{\rm phase0}(k)+|P_1(k)|\right).
\]
When explicit Phase-0 controller telemetry exists, \(N_{\rm phase0}\) is its
\texttt{records\_with\_group\_keys}.  Otherwise the reconstruction uses the
selected-logical screening count recorded by the source payload.
\end{primerbox}

\section*{4. Paper I Table Formulas}

For the non-SNAKE adaptive comparators, the table row uses the explicit
algorithmic ledger
\[
\boxed{
\begin{aligned}
G_K&=\sum_{k=0}^{K}{\tt candidate\_count\_scored}_k,\\
M_K&=\sum_{k=0}^{K}{\tt selector\_metric\_probe\_count}_k,\\
F_K&=\sum_{k=0}^{K}{\tt optimizer\_nfev}_k,\\
S_K&=G_K+M_K+F_K.
\end{aligned}}
\]
For append ADAPT, TETRIS, and QEB, \(M_K=0\) in the Hubbard weak table-facing
sources.  For Geo-ADAPT, \(M_K\) is the geometric selector metric triangle.

For SNAKE, the current Paper I table cell uses the native controller
\texttt{shots\_new} ledger.  Let \(\Wcal_K\) be the controller work scopes
whose \texttt{depth} is in the displayed accepted prefix \(0,\ldots,K\).  Each
scope \(w\) has a phase \(\phi(w)\) and a fresh controller shot count
\[
  c_w={\tt shots\_new}(w).
\]
The reported SNAKE table value is
\[
\boxed{
\begin{aligned}
\Ngrad^{\rm table}(K)
&=
\sum_{\substack{w\in\Wcal_K\\ \phi(w)\in\{0,1\}}} c_w,\\
\Nmetric^{\rm table}(K)
&=
\sum_{\substack{w\in\Wcal_K\\ \phi(w)\in\{2,3\}}} c_w,\\
\Ssnake^{\rm table}(K)
&=
\Ngrad^{\rm table}(K)+\Nmetric^{\rm table}(K).
\end{aligned}
}
\]
Optimizer/refit counts are present in SNAKE telemetry as optimizer provenance;
the current Paper I SNAKE shot-proxy cell is the controller
\texttt{shots\_new} sum.

The SNAKE source still stores Hamiltonian objective telemetry for audits:
\[
  \begin{aligned}
  \NHrefit
  &=
  N_{\rm hist}+N_{\rm resume}+N_{\rm final},\\
  N_{\rm hist}
  &=\sum {\tt history[*].nfev\_opt},\\
  N_{\rm resume}
  &={\tt resume\_boundary\_refit.nfev},\\
  N_{\rm final}
  &={\tt final\_full\_refit.nfev}.
  \end{aligned}
\]
\[
  \NHouter={\tt nfev\_total}-\NHrefit.
\]

\begin{primerbox}{Reported table convention}
For the current Paper I SNAKE row, the promoted support file labels the table
cell as \(\Salg=98\), and the source is \texttt{controller\_shot\_proxy}.
That value is the fresh controller \texttt{shots\_new} total across Phase 0,
Phase I, Phase II, and Phase III.  Reused controller groups appear in the
source as \texttt{shots\_reused}; they are not added again to the table cell.
\end{primerbox}

\section*{5. Actual Reporting Code}

\textbf{\footnotesize A phase count is a measurement-bearing record count.}

\footnotesize\path{pipelines/static_adapt/adapt_pipeline.py}
\begin{lstlisting}
def _controller_work_group_keys_from_records(records):
    records_list = [dict(rec) for rec in (records or []) if isinstance(rec, Mapping)]
    group_keys = []
    records_with_group_keys = 0
    for rec in records_list:
        candidate_term = rec.get("candidate_term")
        if not isinstance(candidate_term, AnsatzTerm):
            idx_raw = rec.get("candidate_pool_index")
            try:
                idx = int(idx_raw)
                if 0 <= idx < len(pool):
                    candidate_term = pool[idx]
            except (TypeError, ValueError):
                candidate_term = None
        if isinstance(candidate_term, AnsatzTerm):
            keys = measurement_group_keys_for_term(candidate_term)
            if keys:
                records_with_group_keys += 1
                group_keys.extend(str(key) for key in keys)
    return group_keys, int(len(records_list)), int(records_with_group_keys)
\end{lstlisting}

\textbf{\footnotesize Phase events that create the reported \(P_j\) counts:
\path{pipelines/static_adapt/adapt_pipeline.py}.}
\begin{lstlisting}
_record_controller_work_for_records(
    controller_work_local,
    snapshot=controller_snapshot_local,
    phase="phase1",
    event_kind=(
        "phase1_insertion_probe"
        if bool(insertion_probe_triggered_local)
        else "phase1_append_probe"
    ),
    records=phase1_records_local,
    depth_value=int(depth + 1),
    scope_qualifiers=branch_scope_qualifiers,
)

phase2_last_shortlist_eval_records_local = [dict(rec) for rec in full_records]
_record_controller_work_for_records(
    controller_work_local,
    snapshot=controller_snapshot_local,
    phase="phase2",
    event_kind="phase2_rerank_records",
    records=full_records,
    depth_value=int(depth + 1),
    scope_qualifiers=branch_scope_qualifiers,
)

if phase3_live_now_local:
    _record_controller_work_for_records(
        controller_work_local,
        snapshot=controller_snapshot_local,
        phase="phase3",
        event_kind="phase3_reduced_geometry_rerank",
        records=phase2_shortlisted_records_local,
        depth_value=int(depth + 1),
        scope_qualifiers=branch_scope_qualifiers,
    )
\end{lstlisting}

\textbf{\footnotesize QWC key construction:
\path{pipelines/scaffold/hh_continuation_scoring.py}.}
\begin{lstlisting}
def measurement_group_specs_for_term(term):
    label_coeffs = _pauli_label_coeffs_from_term(term)
    specs = []
    for label, coeff in label_coeffs:
        label_s = str(label)
        if _pauli_weight_exyz(label_s) <= 0:
            continue
        coeff_abs = abs(complex(coeff))
        if not math.isfinite(float(coeff_abs)) or float(coeff_abs) == 0.0:
            continue
        specs.append(
            MeasurementGroupSpec(
                group_key=label_s,
                coeff_l2=float(coeff_abs),
                term_count=1,
                source="polynomial_coeff_l2_v1",
            )
        )
    if specs:
        return _compress_measurement_group_specs(specs)
    return [
        MeasurementGroupSpec(group_key=str(key), coeff_l2=1.0, term_count=1, source="label_fallback_l2_v1")
        for key in _measurement_group_keys_from_labels(_pauli_labels_from_term(term))
    ]
\end{lstlisting}

\textbf{\footnotesize Runtime reconstruction:
\path{pipelines/exact_bench/snake_table_i_measurement_work.py}.}
\begin{lstlisting}
def _phase_records_with_group_keys(phase_map, phase):
    entry = phase_map.get(phase)
    if entry is None:
        return 0.0, {"phase": phase, "status": "absent_zero"}
    if not isinstance(entry, Mapping):
        return None, {"phase": phase, "status": "invalid_phase_payload"}
    value = _finite_nonnegative(entry.get("records_with_group_keys"))
    if value is None:
        return None, {"phase": phase, "status": "invalid_records_with_group_keys"}
    return value, {"phase": phase, "status": "ok", "records_with_group_keys": float(value)}
\end{lstlisting}

\begin{lstlisting}
phase0, phase0_accounting = _selected_logical_phase0_gradient_work(source_payload, phase_map)
phase1, detail = _phase_records_with_group_keys(phase_map, "phase1")
phase2, detail = _phase_records_with_group_keys(phase_map, "phase2")
phase3, detail = _phase_records_with_group_keys(phase_map, "phase3")

n_h_refit = float(history_refit + resume_refit + final_refit)
n_h_outer = float(nfev_total - n_h_refit)

components = {
    "N_H_outer_eval": float(n_h_outer),
    "N_grad": float(phase0 + phase1),
    "N_metric": float(phase2 + phase3),
    "N_H_refit_eval": float(n_h_refit),
}
\end{lstlisting}

\textbf{\footnotesize Current SNAKE table cell:
\path{pipelines/exact_bench/snake_table_i_measurement_work.py}.}
\begin{lstlisting}
phase1_shots, detail = _phase_shots_new(phase_map, "phase1")
phase2_shots, detail = _phase_shots_new(phase_map, "phase2")
phase3_shots, detail = _phase_shots_new(phase_map, "phase3")

n_grad_shots = float(phase0_shots + phase1_shots)
n_metric_shots = float(phase2_shots + phase3_shots)
total = float(n_grad_shots + n_metric_shots)
meta.update(
    status="ok",
    controller_shot_proxy=total,
    N_grad_controller=float(n_grad_shots),
    N_metric_controller=float(n_metric_shots),
    controller_phase_shots_new={
        "phase0": float(phase0_shots),
        "phase1": float(phase1_shots),
        "phase2": float(phase2_shots),
        "phase3": float(phase3_shots),
    },
)
\end{lstlisting}

\textbf{\footnotesize Non-SNAKE paper-facing total:
\path{pipelines/exact_bench/generic_static_metric_enrichment.py}.}
\begin{lstlisting}
s_alg = (
    weights["s_H_outer"] * components["N_H_outer_eval"]
    + weights["s_grad"] * components["N_grad_probe"]
    + weights["s_metric"] * components["N_metric_probe"]
    + weights["s_H_refit"] * components["N_H_refit_eval"]
)
row_updates = {
    "S_alg": float(s_alg),
    "S_alg_N_H_outer_eval": float(components["N_H_outer_eval"]),
    "S_alg_N_grad_probe": float(components["N_grad_probe"]),
    "S_alg_N_metric_probe": float(components["N_metric_probe"]),
    "S_alg_N_H_refit_eval": float(components["N_H_refit_eval"]),
    "S_alg_N_other_quantum": 0.0,
}
\end{lstlisting}

\section*{6. Hubbard Weak Prefix Examples}

These values are computed from the visible Paper I Hubbard weak sources
(\(\mathrm{Hubbard}\ L=2,\ U/t=0.5\)).  For the comparator rows,
\[
  S_K=G_K+M_K+F_K.
\]
For SNAKE, the current table source contains two accepted rounds; rows
\(K=2,\ldots,5\) are undefined for that table-facing artifact.

\begin{table*}[t]
\centering
\scriptsize
\begin{tabular}{@{}lrlrrrr@{}}
\hline
algorithm & \(K\) & increment & \(G_K\) & \(M_K\) & \(F_K\) & \(S_K\)\\
\hline
Append & 0 & \(18+9037\) & 18 & 0 & 9037 & 9055\\
Append & 1 & \(17+5668\) & 35 & 0 & 14705 & 14740\\
Append & 2 & \(17+1291\) & 52 & 0 & 15996 & 16048\\
Append & 3 & \(17+22\) & 69 & 0 & 16018 & 16087\\
Append & 4 & \(17+22\) & 86 & 0 & 16040 & 16126\\
Append & 5 & \(17+7\) & 103 & 0 & 16047 & 16150\\
\hline
TETRIS & 0 & \(18+16\) & 18 & 0 & 16 & 34\\
TETRIS & 1 & \(17+1759\) & 35 & 0 & 1775 & 1810\\
TETRIS & 2 & \(17+892\) & 52 & 0 & 2667 & 2719\\
TETRIS & 3 & \(17+133\) & 69 & 0 & 2800 & 2869\\
TETRIS & 4 & \(17+154\) & 86 & 0 & 2954 & 3040\\
TETRIS & 5 & \(17+145\) & 103 & 0 & 3099 & 3202\\
\hline
Geo & 0 & \(18+171+9037\) & 18 & 171 & 9037 & 9226\\
Geo & 1 & \(17+153+5668\) & 35 & 324 & 14705 & 15064\\
Geo & 2 & \(17+153+364\) & 52 & 477 & 15069 & 15598\\
Geo & 3 & \(17+153+133\) & 69 & 630 & 15202 & 15901\\
Geo & 4 & \(17+153+16\) & 86 & 783 & 15218 & 16087\\
Geo & 5 & \(17+153+25\) & 103 & 936 & 15243 & 16282\\
\hline
QEB & 0 & \(9+9037\) & 9 & 0 & 9037 & 9046\\
QEB & 1 & \(8+5668\) & 17 & 0 & 14705 & 14722\\
QEB & 2 & \(8+745\) & 25 & 0 & 15450 & 15475\\
QEB & 3 & \(8+214\) & 33 & 0 & 15664 & 15697\\
QEB & 4 & \(8+262\) & 41 & 0 & 15926 & 15967\\
QEB & 5 & \(8+22\) & 49 & 0 & 15948 & 15997\\
\hline
SNAKE & 0 & \(46\) & 22 & 24 & 0 & 46\\
SNAKE & 1 & \(52\) & 44 & 54 & 0 & 98\\
\hline
\end{tabular}
\caption{\footnotesize Paper I Hubbard weak prefix shot-proxy examples.  \(G\)
is gradient/controller-gradient work, \(M\) is metric/controller-metric work,
and \(F\) is refit optimizer work for non-SNAKE comparator rows.}
\end{table*}

For SNAKE the two increments are the native controller event sums
\[
\begin{array}{c|ccccc|c}
K & \phi0 & \phi1 & \phi2 & \phi3_{\rm union} & \phi3_{\rm geom} & \Delta S\\
\hline
0 & 11 & 11 & 11 & 2 & 11 & 46\\
1 & 11 & 11 & 11 & 8 & 11 & 52
\end{array}
\]
The depth-\(K=1\) Phase-0 scope also reports \(11\) reused shots; the table adds
the fresh \texttt{shots\_new}=11.

\section*{7. Implementation Anchors}

\begin{itemize}
\item \footnotesize\path{MATH/paper_details/static_adapt_paper_I.tex}
\item \footnotesize\path{pipelines/static_adapt/adapt_pipeline.py}
\item \footnotesize\path{pipelines/static_adapt/selector_measurement_proxy.py}
\item \footnotesize\path{pipelines/scaffold/hh_continuation_scoring.py}
\item \footnotesize\path{pipelines/exact_bench/snake_table_i_measurement_work.py}
\item \footnotesize\path{pipelines/exact_bench/generic_static_metric_enrichment.py}
\item \footnotesize\path{MATH/paper_facing/paper_I_static_scaffold/paper_i_tables_i_ii_repeat_enabled_comparator_promotion_20260610.json}
\end{itemize}

\begin{primerbox}{Memory hook}
For non-SNAKE Paper I reporting, the table adds gradient probes, metric probes,
and optimizer refit evaluations.  For the current SNAKE Paper I row, the table
adds fresh controller \texttt{shots\_new} from Phase 0, Phase I, Phase II, and
Phase III; reused controller groups are recorded but not added again.
\end{primerbox}

\end{document}
